\newpage
\phantomsection
\label{prf:archimedean-property}

\begin{remark*}[Return]
\hyperref[thm:archimedean-property]{Return to Theorem}
\end{remark*}

\begin{theorem*}[Archimedean Property]
For every $x \in \mathbb{R}$, there exists $n \in \mathbb{N}$ such that $n > x$.
\end{theorem*}

\begin{proof}
\textbf{Professional Standard Proof.}~\\
Assume for sake of contradiction that the property is false. Then there exists $x \in \mathbb{R}$ such that $n \le x$ for all $n \in \mathbb{N}$. This implies $\mathbb{N}$ is bounded above. By the Least Upper Bound Property of $\mathbb{R}$, $s = \sup \mathbb{N}$ exists. Since $s$ is the least upper bound, $s-1$ is not an upper bound. Thus, there exists $m \in \mathbb{N}$ such that $m > s-1$, which implies $m+1 > s$. Since $m+1 \in \mathbb{N}$, this contradicts that $s$ is an upper bound for $\mathbb{N}$. Hence, $\mathbb{N}$ is unbounded.
\end{proof}

\begin{proof}
\textbf{Detailed Learning Proof.}~\\

\textbf{Step 1. Assume the negation for contradiction.}
We assume that the set of natural numbers $\mathbb{N}$ is bounded above. Formally:
\[ \exists x \in \mathbb{R} \; \forall n \in \mathbb{N} : n \le x \]

\textbf{Step 2. Invoke the Least Upper Bound Property.}
Since $\mathbb{N}$ is non-empty ($1 \in \mathbb{N}$) and assumed to be bounded above, the completeness of $\mathbb{R}$ ensures that $s = \sup \mathbb{N}$ exists in $\mathbb{R}$.

\textbf{Step 3. Apply the $\varepsilon$-characterization of the supremum.}
Let $\varepsilon = 1$. By the characterization of the supremum, there must be some element of the set strictly greater than $s - \varepsilon$. Thus, there exists $m \in \mathbb{N}$ such that:
\[ m > s - 1 \]

\textbf{Step 4. Construct a contradiction within the set.}
Adding $1$ to both sides of the inequality from Step 3:
\[ m + 1 > s \]
By the inductive property of natural numbers, if $m \in \mathbb{N}$, then $m + 1 \in \mathbb{N}$. However, we have found a natural number $m+1$ that is strictly greater than $s$, the supposed upper bound of $\mathbb{N}$.

\textbf{Step 5. Conclude the proof.}
The existence of $m+1 > s$ contradicts the definition of $s$ as an upper bound. Therefore, our initial assumption that $\mathbb{N}$ is bounded must be false. Thus, for every $x \in \mathbb{R}$, there exists some $n \in \mathbb{N}$ such that $n > x$.
\end{proof}

\begin{remark*}[Proof structure]
The proof utilizes a contradiction strategy based on the completeness of the real numbers. It shows that if $\mathbb{N}$ were bounded, its "ceiling" would have to be surpassable by simply adding $1$ to a value near that ceiling.
\end{remark*}

\begin{remark*}[Dependencies]~\\
\begin{itemize}
  \item \hyperref[pred:least-upper-bound-property]{Least Upper Bound Property}
  \item \hyperref[pred:supremum]{Definition of Supremum}
  \item \hyperref[pred:upper-bound]{Definition of Upper Bound}
\end{itemize}
\end{remark*}